


We now share details on our experimental setup.

\noindent\textbf{Input problems.}
We randomly sample 10 problems for each difficulty level from 1 to 5 from the FormalGeo-7K dataset, resulting in a total of 50 problems. Levels correspond to the number of steps in the ground-truth proof. % (from one step at level 1, to five steps at level 5).

\noindent\textbf{Base models.} 
We evaluate \emph{reasoning-capable} models from multiple families using recommended configurations. 
We evaluate OpenAI's o1~\cite{jaech2024openai}, GPT-5~\cite{gpt-5}, Gemini-2.5-Flash~\cite{comanici2025gemini} and Claude Sonnet 4.6~\cite{claude_sonnet_4.6}. Throughout the paper, \emph{baseline} refers to the non-augmented variant of each of these models (i.e., without our neuro-symbolic method), rather than to any specific model.


We use OpenAI o1 with the default settings; GPT-5 with medium \emph{reasoning}; Gemini-2.5-Flash with \emph{t=0.7}; and Claude Sonnet 4.6 with \emph{t=1.0}. Full configuration details are provided in our codebase.


We additionally evaluated \emph{non-reasoning} models, including GPT-4 and GPT-4o~\cite{achiam2023gpt, hurst2024gpt}, and found that they frequently generate proofs with numerous syntax errors. 
Smaller open-weight models such as Qwen-3-8B and Qwen-3-14B~\cite{yang2025qwen3} similarly struggle to adhere to the required output format.






\noindent\textbf{Few-shot variants.}  
For each base model, we evaluate two variants: an analogy-based variant using our method, and a baseline variant without it.
\begin{compactitem}
  \item \textbf{Analogy-based:} We provide $k$ similar (analogous) problems and their proofs as few-shot examples, along with an abridged theorem dictionary which includes only theorems from the proofs of similar problems (Section~\ref{subsec:llm_solver}).
  \item \textbf{Base model (non-analogy):} We provide $k$ random problems and their proofs as few-shot examples, along with the complete dictionary.
\end{compactitem}

 
% \onote{where to mention abridged GDL is also an advantage of our method?}

% \noindent {\bf A note on baselines.}
% This work aims to evaluate how our method (combining analogical guidance and symbolic verification) can contribute to \emph{general-purpose} LLMs in formal proof generation. Accordingly, we do not compare against specialized geometry solvers, as our focus is on augmenting existing LLMs rather than competing with domain-specific systems (see Section~\ref{sec:related_work}).


\noindent {\bf A note about $k$.} In our preliminary exploration, no model was able to correctly follow the proof format with zero-shot prompts (unsurprisingly). We have found that $k=5$ works well in practice; testing the effect of $k$ in more depth is left for future work.





\noindent {\bf Feedback iterations and multiple runs.}
For each base model and variant (analogy-based and non-analogy), we evaluate several inference settings involving verifier-guided feedback and repeated runs.
As outlined in Section~\ref{subsec:verifier}, we integrate the LLM into a \emph{feedback loop} guided by a verifier, allowing the LLM up to five iterations (retries following feedback) per run.  
In addition, since LLMs sometimes fail to progress toward the correct solution (even with feedback), we allow restarting from scratch up to three times per problem.
Within each run, the conversation is append-only: failed attempts and verifier feedback are retained, and each retry conditions on the full interaction history within that run. The $n$ runs are independent, each starting from a fresh context.

% (i.e., three separate runs, with the default -- and unchangeable -- temperature parameter set to 1).


Overall, we evaluate four base models, each under two prompting variants and three inference settings: single-pass, verifier-guided feedback iterations, and multiple independent runs with feedback iterations. 
In total, this yields 24 configurations.

We denote the different configurations by $\text{base}_{i,j}$ or $\text{ours}_{i,j}$,
where $i$ is the number of runs and $j$ the number of verifier-guided retries.


\noindent {\bf Metric.}
We report the fraction of problems for which the model produces a correct proof, for the following settings:
\begin{compactitem}
    \item \textbf{First run, no retries (\%):} The model produces a correct proof on its first attempt.
    \item \textbf{First run, with retries (\%):} Any of $m$ retries (following feedback) in the first run succeeds.
    \item \textbf{Multiple runs, with retries (\%):}  Any attempt across $n$ independent runs (each with up to $m$ retries) succeeds.
\end{compactitem}





More runs and retries increase costs; we found $m=5$, $n=3$ to provide a good tradeoff between performance and budget, and analyze their effect in the next section.
Based on logged API calls, token usage, and current pricing, evaluating 50 problems costs approximately \$100 (o1), \$3 (GPT-5), \$1 (Gemini-2.5-Flash), and \$10 (Claude Sonnet 4.6). The baseline is approximately 3$\times$ more expensive, since it prompts with the full $\sim$18K-token theorem dictionary rather than our $\sim$2.5K-token analogy-derived dictionary, resulting in substantially larger prompts and typically more verifier-guided retries.



% We also measure the average number of runs and retries, and analyze the distribution of errors in the different tiers (see Section~\ref{subsec:verifier_error_analysis}).
